   \documentclass[instructions.tex]{subfiles}
\begin{document}

\chapter{On risque encore plus quand on remet sa conversion au temps de la maladie.}

Comprenez-vous à quoi vous vous exposez, lorsque vous remettez votre conversion au temps d’une maladie ? pourrez-vous alors vous convertir ? et, en supposant que vous le puissiez, le ferez-vous ?

\primo La maladie ne vous en laissera peut-être ni le pouvoir ni le temps. L’ardeur d’une fièvre et la langueur vous affaibliront de telle sorte, qu’il ne vous sera presque plus possible de régler votre conscience. Dieu veuille que vous ne l’éprouviez pas, comme tant d’autres l’ont éprouvé.

Un malade qu’on exhortait à produire des actes de religion : \textit{Je ne le puis}, répondit-il : \textit{j’ai l’esprit trop faible ; si j’avais attendu jusqu’à présent de me préparer à la mort, mon salut serait désespéré}. Un religieux, d’une éminente vertu, disait, après une longue maladie : \textit{Lorsqu’on tâchait de m’inspirer quelques sentiments de Dieu, j’avais l’esprit si faible, que je n’entendais que des paroles confuses, sans en comprendre le sens}. S’il en est ainsi des gens de bien, aurez-vous plus de facilité qu’eux, vous qu’une longue habitude dans le vice et dans la dissipation, aura rendu presque incapable de faire un retour sur vous-même ?

Si vous mourez sans amour de Dieu, dans l’attache au péché mortel, votre damnation sera inévitable. Croyez-vous qu’étant malade vous apprendrez dans un quart d’heure à aimer Dieu souverainement, à détacher votre cœur du péché et des créatures, à mettre ordre aux embarras de votre conscience ? Il faudra pour cela des réflexions très-sérieuses : serez-vous en état de les faire ? Il faudra que votre cœur et votre esprit soient changés ; mais le cœur et l’esprit de l’homme se changent-ils aussi facilement ? Vous avez de la peine à le faire en santé, et vous prétendez le faire étant malade ! Oui, vous prétendez apprendre la plus difficile de toutes les sciences, qui est la science de bien mourir, dans un temps où vous ne serez presque plus capable de rien ! Où est votre raison ? Si vous vouliez apprendre à écrire ou à dessiner, attendriez-vous le temps d’une maladie mortelle ? pourquoi donc attendez-vous le temps de la maladie pour vous disposer à une mort sainte ? Il est bien plus difficile de mourir en pénitent et en saint, que d’être écrivain ou architecte.

Prétendre qu’à sa mort on se convertira et qu’on apprendra à mourir en saint, après avoir toujours vécu en pécheur ; prétendre qu’on mourra dans la ferveur, après avoir toujours vécu dans la lâcheté et l’oubli de Dieu, c’est imiter un téméraire qui prétendrait parler tout-à-coup une langue qu’il n’aurait jamais apprise. Oh ! qu’il est bien tard d’apprendre à bien vivre, quand on est sur le point de mourir ! Qu’il est difficile alors de penser en saint, après avoir toujours pensé en pécheur ! Qu’il est bien tard de commencer, quand il faut finir !

\secundo Vous pourriez alors vous convertir par une grâce extraordinaire ; mais avez-vous lieu de vous attendre à une telle grâce, en continuant de vivre dans le crime ? Et en supposant que vous aurez la grâce et le temps de vous convertir dans une maladie, en aurez-vous la volonté ? cette volonté sera-t-elle sincère ? ne sera-t-elle point une volonté imaginaire et sans effet, qui ne changera point votre cœur ? Vous devez d’autant plus le craindre, que saint Augustin a dit que la conversion tardive est souvent nulle, et que la conversion des mourans est ordinairement morte. Pourquoi ? parce qu’on ne quitte pas alors le péché ; c’est plutôt le péché qui nous quitte ; et qu’on ne cesse de vouloir pécher, que parce qu’on ne peut plus pécher\footnote{Peccata tua te dimiserunt, non tu illa.}.

Une telle conversion a bien du rapport avec celle des criminels, qui déclarent tout quand ils sont à la torture, et qui désavouent un moment après tout ce qu’ils ont déclaré.

Pharaon fut puni de Dieu plusieurs fois de suite : à chaque fois il promettait d’obéir à Dieu ; mais dès qu’il ne sentait plus les châtimens du Tout-Puissant, il oubliait ses promesses. La même chose arrive à la plupart de ceux qui diffèrent leur conversion jusqu’à l’extrémité, ils font une confession et des promesses qu’ils ne feraient point s’ils n’étaient pas en danger de mourir. Ils sont épouvantés par les horreurs de la mort ; mais sont-ils convertis ? Il est aisé d’en juger par le peu d’amendement, et par leur rechute après la maladie.

C’est une grande folie de négliger un vent favorable et d’attendre la tempête pour se mettre en mer. Mais c’est encore une plus grande folie, dit saint Eucher, d’attendre au moment le plus incertain de la vie pour assurer son salut. Est-ce en effet une conduite prudente de fonder son salut sur quelques propos de mieux vivre, lorsqu’on n’a plus de temps à vivre ? Pensez-y sérieusement ; pensez-y effectivement : c’est chercher bien tard le remède, quand le mal est extrême.

\section{Résolutions}

\primo Je me défierai toujours de la maladie, comme d’un temps peu propre à l’exécution des grandes résolutions qu’on n’a pas eu le courage de mettre en pratique pendant qu’on jouissait de la santé.

\secundo Et je regarderai le retardement de la conversion au temps de la maladie, comme une injure faite à Dieu, bien capable d’attirer sa colère, et d’éloigner ses grâces de moi.

 Mon Dieu, c’en est fait, je reviens à vous : recevez-moi comme le père de famille reçut l’enfant prodigue de l’Évangile.

\end{document}
